% GENERATED FILE. Edit canonical lessons, then run npm run generate:books.
% canonical-source-hash: fnv1a64:a032794a0bba8284
% canonical-lessons: FR-C364-justice, FR-C364-injuste, FR-C364-injustice, FR-C364-prejuge, FR-C364-stereotype

\chapter{Juste et injuste}
\label{ch:fr-juste-et-injuste}
\hlchaptermodality{364}

\begin{chapteropening}
\textbf{By the end of this chapter:} \emph{I can talk about justice and injustice, say what is unfair, and name prejudice and stereotypes.}

Complete the last lesson of chapter 364: I can talk about justice and injustice, say what is unfair, and name prejudice and stereotypes.
\end{chapteropening}

\section[la justice]{la justice — the justice; the courts}

\label{lesson:FR-C364-justice}



\begin{quote}
Before the new one: say the French for the banner, then the French for to protest.
\end{quote}

\subsection*{You'll want to know: la justice}
\textbf{la justice} — \textquotedblleft{}the justice; the courts\textquotedblright{}.

\begin{quote}
La famille demande justice.
\end{quote}

\begin{quote}
\emph{The family is asking for justice.}
\end{quote}

\emph{La justice} is also the legal system: \emph{aller en justice}.

\begin{cousinweb}
\emph{justice} comes through Old French from Latin \emph{iūstitia} \textquotedblleft{}fairness, justice\textquotedblright{} (EtymDB), from \emph{iūstus} \textquotedblleft{}right, fair\textquotedblright{} (Lewis \& Short).
\end{cousinweb}

\begin{grammarlens}[title={What you've built}]
One more word for telling, explaining and describing at B1.
\end{grammarlens}

\subsection*{Your turn}
\begin{itemize}
\raggedright
  \item \emph{Say it:} \emph{la justice}
  \item \emph{Say it:} \emph{la justice}, once more
  \item \emph{Say it:} a sentence of your own with \emph{la justice}
  \item \emph{Recall:} say the French for the banner, then the French for to protest, then say \emph{la justice} again
\end{itemize}

\begin{tcolorbox}[breakable,colback=teal!4,colframe=teal!35!black,title={Before you move on}]
What does la justice mean? (\textquotedblleft{}The justice; the courts\textquotedblright{}.) Say it once more.
\end{tcolorbox}
\section[injuste]{injuste — unfair, unjust}

\label{lesson:FR-C364-injuste}



\begin{quote}
Before the new one: say the French for to protest, then the French for the justice.
\end{quote}

\subsection*{You'll want to know: injuste}
\textbf{injuste} — \textquotedblleft{}unfair, unjust\textquotedblright{}.

\begin{quote}
C'est injuste : il a eu la même note que moi sans travailler.
\end{quote}

\begin{quote}
\emph{It is unfair: he got the same mark as me without working.}
\end{quote}

One form for both genders.

\begin{cousinweb}
\emph{injuste} is a learned borrowing from Latin \emph{iniūstus} \textquotedblleft{}unjust\textquotedblright{} (EtymDB; Lewis \& Short), built from \emph{in-} \textquotedblleft{}not\textquotedblright{} and \emph{iūstus} \textquotedblleft{}right, fair\textquotedblright{}.
\end{cousinweb}

\begin{grammarlens}[title={What you've built}]
One more word for telling, explaining and describing at B1.
\end{grammarlens}

\subsection*{Your turn}
\begin{itemize}
\raggedright
  \item \emph{Say it:} \emph{injuste}
  \item \emph{Say it:} \emph{injuste}, once more
  \item \emph{Say it:} a sentence of your own with \emph{injuste}
  \item \emph{Recall:} say the French for to protest, then the French for the justice, then say \emph{injuste} again
\end{itemize}

\begin{tcolorbox}[breakable,colback=teal!4,colframe=teal!35!black,title={Before you move on}]
What does injuste mean? (\textquotedblleft{}Unfair, unjust\textquotedblright{}.) Say it once more.
\end{tcolorbox}
\section[l'injustice]{l'injustice — the injustice}

\label{lesson:FR-C364-injustice}



\begin{quote}
Before the new one: say the French for the justice, then the French for unfair.
\end{quote}

\subsection*{You'll want to know: l'injustice}
\textbf{l'injustice} — \textquotedblleft{}the injustice\textquotedblright{}.

\begin{quote}
L'injustice la met en colère.
\end{quote}

\begin{quote}
\emph{Injustice makes her angry.}
\end{quote}

Feminine: \emph{une injustice}.

\begin{cousinweb}
\emph{injustice} is a learned borrowing from Latin \emph{iniūstitia} \textquotedblleft{}injustice\textquotedblright{} (Lewis \& Short), from \emph{iniūstus} \textquotedblleft{}unjust\textquotedblright{}, built on \emph{iūstus} \textquotedblleft{}right, fair\textquotedblright{}.
\end{cousinweb}

\begin{grammarlens}[title={What you've built}]
One more word for telling, explaining and describing at B1.
\end{grammarlens}

\subsection*{Your turn}
\begin{itemize}
\raggedright
  \item \emph{Say it:} \emph{l'injustice}
  \item \emph{Say it:} \emph{l'injustice}, once more
  \item \emph{Say it:} a sentence of your own with \emph{l'injustice}
  \item \emph{Recall:} say the French for the justice, then the French for unfair, then say \emph{l'injustice} again
\end{itemize}

\begin{tcolorbox}[breakable,colback=teal!4,colframe=teal!35!black,title={Before you move on}]
What does l'injustice mean? (\textquotedblleft{}The injustice\textquotedblright{}.) Say it once more.
\end{tcolorbox}
\section[le préjugé]{le préjugé — the prejudice}

\label{lesson:FR-C364-prejuge}



\begin{quote}
Before the new one: say the French for unfair, then the French for the injustice.
\end{quote}

\subsection*{You'll want to know: le préjugé}
\textbf{le préjugé} — \textquotedblleft{}the prejudice\textquotedblright{}.

\begin{quote}
Elle a beaucoup de préjugés sur les jeunes.
\end{quote}

\begin{quote}
\emph{She has a lot of prejudices about young people.}
\end{quote}

Literally judged in advance: \emph{pré-} + \emph{jugé}.

\begin{cousinweb}
\emph{préjugé} is built from \emph{pré-} \textquotedblleft{}before\textquotedblright{} and \emph{jugé} \textquotedblleft{}judged\textquotedblright{}, after Latin \emph{praeiūdicium} \textquotedblleft{}a judgement made beforehand\textquotedblright{} (Lewis \& Short), built on \emph{iūdex} \textquotedblleft{}judge\textquotedblright{}.
\end{cousinweb}

\begin{grammarlens}[title={What you've built}]
One more word for telling, explaining and describing at B1.
\end{grammarlens}

\subsection*{Your turn}
\begin{itemize}
\raggedright
  \item \emph{Say it:} \emph{le préjugé}
  \item \emph{Say it:} \emph{le préjugé}, once more
  \item \emph{Say it:} a sentence of your own with \emph{le préjugé}
  \item \emph{Recall:} say the French for unfair, then the French for the injustice, then say \emph{le préjugé} again
\end{itemize}

\begin{tcolorbox}[breakable,colback=teal!4,colframe=teal!35!black,title={Before you move on}]
What does le préjugé mean? (\textquotedblleft{}The prejudice\textquotedblright{}.) Say it once more.
\end{tcolorbox}
\section[le stéréotype]{le stéréotype — the stereotype}

\label{lesson:FR-C364-stereotype}



\begin{quote}
Before the new one: say the French for the injustice, then the French for the prejudice.
\end{quote}

\subsection*{You'll want to know: le stéréotype}
\textbf{le stéréotype} — \textquotedblleft{}the stereotype\textquotedblright{}.

\begin{quote}
Ce film est plein de stéréotypes.
\end{quote}

\begin{quote}
\emph{This film is full of stereotypes.}
\end{quote}

First a word from printing: a solid metal plate that copies a page.

\begin{cousinweb}
\emph{stéréotype} is usually explained as a French formation from Greek \emph{stereós} \textquotedblleft{}solid\textquotedblright{} and \emph{týpos} \textquotedblleft{}mark, impression\textquotedblright{}: first a printing plate, then a fixed idea that is copied without thought.
\end{cousinweb}

\begin{grammarlens}[title={What you've built}]
One more word for telling, explaining and describing at B1.
\end{grammarlens}

\subsection*{Your turn}
\begin{itemize}
\raggedright
  \item \emph{Say it:} \emph{le stéréotype}
  \item \emph{Say it:} \emph{le stéréotype}, once more
  \item \emph{Say it:} a sentence of your own with \emph{le stéréotype}
  \item \emph{Recall:} say the French for the injustice, then the French for the prejudice, then say \emph{le stéréotype} again
\end{itemize}

\begin{tcolorbox}[breakable,colback=teal!4,colframe=teal!35!black,title={Before you move on}]
What does le stéréotype mean? (\textquotedblleft{}The stereotype\textquotedblright{}.) Say it once more.
\end{tcolorbox}
